%----------------------------------------------------------------------------------------
% TAILORED CV — University of Toronto
% Wheeler Microfluidics Laboratory — Digital Microfluidics & EWOD
% PI: Prof. Aaron Wheeler | aaron.wheeler@utoronto.ca
% Deadline: 31 May 2026 (BME Round 2)
%----------------------------------------------------------------------------------------

\documentclass[11pt,a4paper]{article}
\usepackage[top=1.8cm,bottom=1.8cm,left=2cm,right=2cm]{geometry}
\usepackage{enumitem}
\usepackage{titlesec}
\usepackage{hyperref}
\usepackage{parskip}
\usepackage{tabularx}
\usepackage{xcolor}
\usepackage{fancyhdr}
\usepackage{lastpage}
\usepackage[T1]{fontenc}
\usepackage[utf8]{inputenc}
\usepackage{textcomp}
\usepackage{microtype}

\definecolor{accent}{RGB}{0,70,127}

\hypersetup{colorlinks=true, urlcolor=accent, linkcolor=accent}

\titleformat{\section}{\large\bfseries\color{accent}}{}{0em}{}[\titlerule]
\titlespacing{\section}{0pt}{10pt}{5pt}

\pagestyle{fancy}
\fancyhf{}
\renewcommand{\headrulewidth}{0pt}
\fancyfoot[C]{\small\color{gray}
  Hamza Abbas $\cdot$ CV $\cdot$ Page \thepage\ of \pageref{LastPage}}

\setlist[itemize]{leftmargin=*, noitemsep, topsep=2pt}

%========================================================================================
\begin{document}
%========================================================================================

%--- HEADER ---
\begin{center}
  {\LARGE\bfseries\color{accent} Hamza Abbas}\\[4pt]
  Lecturer, Faculty of Mechanical Engineering $\cdot$ GIK Institute, Topi, KPK,
  Pakistan\\[3pt]
  \href{mailto:hamza.abbas0133@gmail.com}{hamza.abbas0133@gmail.com}
  \quad$|$\quad +92\,348\,0555875
  \quad$|$\quad
  \href{https://linkedin.com/in/hamza-abbas1}{linkedin.com/in/hamza-abbas1}
\end{center}

%--- PROFILE ---
\section{Profile}

Mechanical engineer and lecturer with hands-on expertise in
\textbf{digital microfluidic (DMF) device fabrication, EWOD actuation, paper-based
diagnostics, and capillary flow modelling}. As a Graduate Research Assistant at GIK
Institute's Biofuels \& Microfluidics Lab, I fabricated and tested paper-based and
digital microfluidic devices (CO$_2$ laser micromachining, photolithographic electrode
patterning on PET, EWOD), and modelled capillary transport using COMSOL Multiphysics.
My modified Lucas--Washburn model predicted wetting-front propagation in heterogeneous
porous substrates to below 8\% error. Two manuscripts from this work are under review
at the \textit{Journal of Fluid Mechanics} and \textit{Langmuir} (ACS). I am
particularly drawn to the Wheeler Lab's work on digital microfluidics and EWOD-driven
point-of-care diagnostics, where my fabrication and capillary-flow skills directly
contribute.

%--- RESEARCH EXPERIENCE (LEAD) ---
\section{Research Experience — Digital Microfluidics \& EWOD}

\textbf{Graduate Research Assistant} -- Biofuels \& Microfluidics Lab, GIK
Institute \hfill \textit{Feb 2022 -- May 2024}\\
\textit{Topi, Pakistan}
\begin{itemize}
  \item \textbf{EWOD device fabrication \& actuation:} Patterned ITO electrodes
        on PET film via photolithography; applied dielectric and hydrophobic
        coatings (SU-8, Cytop); actuated droplets via EWOD and characterised
        contact-angle response as a function of applied voltage and coating
        selection (published, \textit{MATEC Web of Conferences}, 2024).
  \item \textbf{Paper-based DMF devices:} Designed and fabricated lateral-flow
        paper microfluidic devices using CO$_2$ laser micromachining and thermal
        lamination; systematically varied reservoir thickness, contact area, and
        fluid volume to control flow velocity (manuscripts under review).
  \item \textbf{Capillary flow modelling:} Derived a modified Lucas--Washburn
        equation for wetting-front propagation in stratified porous assemblies;
        validated to below 8\% error — directly applicable to passive
        capillary-driven DMF systems.
  \item \textbf{COMSOL simulation:} Modelled electric-field distributions,
        porous-media flow, and heat transfer in microfluidic geometries; compared
        computational outputs against experimental data.
  \item \textbf{Imaging \& analysis:} Tracked wetting fronts and droplet
        actuation using ImageJ; performed curve fitting and statistical analysis
        using OriginPro and MATLAB.
  \item \textbf{Complex-geometry flow:} Investigated resistive fluid transport
        in a hierarchical leaf-inspired channel architecture (under review,
        \textit{Journal of Fluid Mechanics}, 2025).
\end{itemize}

\vspace{4pt}
\textbf{Project Research Assistant} -- HEC Project TTSF-74, GIK Institute
\hfill \textit{Feb 2022 -- Feb 2024}\\
\textit{Topi, Pakistan}
\begin{itemize}
  \item Developed a flexible capillary-driven time-temperature indicator on
        porous substrates for food quality monitoring; supported commercialisation
        through spin-off startup ``Flow Time.''
\end{itemize}

%--- PUBLICATIONS ---
\section{Publications}

\textit{Journal Articles -- Under Review}
\begin{itemize}
  \item Ali, M.; Jafry, A.T.; \textbf{Abbas, H.}; Lee, M.; Abbasi, M.S.; Shah,
        S.F.; Ajab, H. ``Revealing the Mechanics of Highly Resistive Fluid Flow in
        a Leaf Architecture using Paper Microfluidics.''
        \textbf{\textit{Journal of Fluid Mechanics}} (Cambridge), under review,
        2025.
  \item \textbf{Abbas, H.}; Jafry, A.T.; Ali, M.; Ajab, H.; Cheema, T.A.
        ``Impact of paper thickness, contact area, and fluid volume on flow
        dynamics in paper-based microfluidics.''
        \textbf{\textit{Langmuir}} (ACS), under review, 2025.
\end{itemize}

\vspace{4pt}
\textit{Conference Proceedings}
\begin{itemize}
  \item \textbf{Abbas, H.}; Ali, M.; Ullah, H.; Jafry, A.T. ``Actuation and Flow
        Control in Paper-Based Microfluidics by Varying Thickness of Storage
        Reservoir.'' \textbf{\textit{MATEC Web of Conferences}}, 2023.
  \item \textbf{Abbas, H.}; Ali, M.; Naeem, N.; Ullah, H.; Jafry, A.T. ``Effect
        of Storage Reservoir on Fluid Velocity in Lateral Flow Paper Device.''
        \textbf{\textit{Engineering Proceedings}} (MDPI), 2023.
  \item Ullah, H.; \textbf{Abbas, H.}; et al. ``Droplet Actuation Enhancement
        through Voltage Control and Hydrophobic Coating Selection.''
        \textbf{\textit{MATEC Web of Conferences}}, 2024.
\end{itemize}

%--- EDUCATION ---
\section{Education}

\textbf{M.S. in Mechanical Engineering} -- Thermofluid Specialisation, GIK
Institute \hfill \textit{Feb 2022 -- Jun 2024}\\
\textit{Topi, Pakistan}
\begin{itemize}
  \item GPA: 3.67/4.00. Merit-Based Graduate Scholarship and Assistantship (GA-1).
  \item Thesis: \textit{Dynamics of Fluid Flow Between Two Porous Membranes.}
        Modified Lucas--Washburn model for stratified porous assemblies with
        dissimilar permeabilities; validated to below 8\% error.
  \item Key courses: Microfluidic Systems (A-), Intermediate Fluid Mechanics (A),
        Mass/Momentum/Energy Transport (A-), Energy Management \& Auditing (A).
\end{itemize}

\vspace{4pt}
\textbf{B.S. in Mechanical Engineering} -- GIK Institute
\hfill \textit{Aug 2017 -- Jul 2021}\\
\textit{Topi, Pakistan}
\begin{itemize}
  \item Senior Design Project: integrated ORC+VARS waste heat recovery system
        modelled in Aspen Plus (17.9\% combined efficiency, 8.7\,MW net output).
\end{itemize}

%--- ACADEMIC POSITION ---
\section{Academic Position}

\textbf{University Lecturer} -- GIK Institute, Faculty of Mechanical Engineering
\hfill \textit{Sep 2024 -- Present}\\
\textit{Topi, Pakistan}
\begin{itemize}
  \item Teach HVAC, Engineering Dynamics, CAD/CAM Lab, and Thermofluids-I lab;
        manage the Fluid Mechanics Laboratory.
  \item LMS Coordinator; OBE documentation aligned with Washington Accord.
  \item Organised the 6th International Workshop on Functional Reverse Engineering
        of Machine Tools (WRE~2024).
\end{itemize}

%--- TECHNICAL SKILLS ---
\section{Technical Skills}

\begin{tabularx}{\textwidth}{@{}>{\bfseries}lX@{}}
DMF / EWOD Fabrication
  & Photolithographic ITO electrode patterning on PET, dielectric \& hydrophobic
    coating (SU-8, Cytop), CO$_2$ laser micromachining, thermal lamination,
    contact-angle characterisation \\[3pt]
Numerical Simulation
  & COMSOL Multiphysics (microfluidic flow, electric fields, porous-media
    transport, heat transfer) \\[3pt]
Analysis \& Coding
  & MATLAB, Python, OriginPro (curve fitting), ImageJ (droplet/wetting-front
    tracking) \\[3pt]
Process Simulation
  & Aspen Plus (ORC, VARS, thermodynamic cycle modelling) \\[3pt]
Other
  & English C1 (CEFR); Microsoft Office; Moodle LMS \\
\end{tabularx}

%--- MEMBERSHIPS & AWARDS ---
\section{Professional Memberships \& Awards}
\begin{itemize}
  \item Registered Mechanical Engineer, Pakistan Engineering Council
        (Lifetime member) -- Sep 2021.
  \item Merit-Based Graduate Scholarship \& Assistantship (GA-1), GIK Institute
        -- Feb 2022.
  \item 5th place, DICE AFS-2022 Innovation Competition, MNS University of
        Agriculture, Multan -- Oct 2022.
\end{itemize}

%--- REFERENCES ---
\section{References}

\begin{tabularx}{\textwidth}{@{}XX@{}}
\textbf{Dr Ali Turab Jafry}       & \textbf{Dr Taqi Ahmad Cheema} \\
Associate Professor, FME          & Dean / Associate Professor, FME \\
GIK Institute, Topi               & GIK Institute, Topi \\
\href{mailto:ali.turab@giki.edu.pk}{ali.turab@giki.edu.pk}
  & \href{mailto:tacheema@giki.edu.pk}{tacheema@giki.edu.pk} \\
\end{tabularx}

\end{document}
